Finite Signals, Windows, and Sinc-Like DTFTs

Geometric sum (r not equal to 1)

$$\sum_{n=0}^{N-1}r^n=\frac{1-r^N}{1-r}$$

Infinite sum ($|r|<1$)

$$\sum_{n=0}^{\infty}r^n=\frac{1}{1-r}$$

Rectangular sequence starting at n=0

$$r_N[n]=1,\quad 0\le n<N$$

$$R_N(e^{j\omega})=e^{-j\omega(N-1)/2}\frac{\sin(N\omega/2)}{\sin(\omega/2)}$$

At $\omega=0$, the limiting value is N.

Symmetric rectangle (N=2M+1)

$$r[n]=1,\quad -M\le n\le M$$

$$D_N(\omega)=\frac{\sin(N\omega/2)}{\sin(\omega/2)}$$

$$R(e^{j\omega})=D_N(\omega)$$

Periodic sinc-like function is also called a Dirichlet kernel.

At $\omega=2\pi m$, use $D_N(\omega)=N$ by continuity.

Finite centered sine, $-M\le n\le M$:

$$x[n]=\sin(\omega_0 n)r[n]$$

$$X(e^{j\omega})=\frac{D_N(\omega-\omega_0)-D_N(\omega+\omega_0)}{2j}$$

Finite centered cosine, $-M\le n\le M$:

$$x[n]=\cos(\omega_0 n)r[n]$$

$$X(e^{j\omega})=\frac{D_N(\omega-\omega_0)+D_N(\omega+\omega_0)}{2}$$

The two shifted kernels center near $\omega=\pm\omega_0$. Truncation spreads the spectrum, and exact magnitude peaks may shift.

Odd finite-sine expansion

$$X(e^{j\omega})=-2j\sum_{n=1}^{M}\sin(\omega_0 n)\sin(\omega n)$$

For a finite sine symmetric about 0:

$$X(e^{j0})=X(e^{j\pi})=0$$

Frequency content:

Dominant: large $|X(e^{j\omega})|$.

Weak: small nonzero $|X(e^{j\omega})|$.

Absent: $X(e^{j\omega})=0$.
